\section{Tensor product of representations}
\lecture\label{lec:7}

\index{Representation!tensor product of}
Let $V$ be a vector space with basis $\{v_1,\dots,v_k\}$ and 
$W$ be a vector space with basis $\{w_1,\dots,w_l\}$. A 
\emph{tensor product} of $V$ and $W$ is a vector space $X$ 
together with a bilinear map 
\[
V\times W\to X,
\quad
(v,w)\mapsto v\otimes w,
\]
such that $\{v_i\otimes w_j:1\leq i\leq k,\,1\leq j\leq l\}$ is a  
basis of $X$. The tensor product of $V$ and $W$ is unique up to isomorphism 
and is denoted by $V\otimes W$. Note that
\[
\dim(V\otimes W)=(\dim V)(\dim W).
\]
Moreover, every element of $V\otimes W$ is a finite sum 
of the form
\[
\sum_{i,j}\lambda_{ij}v_i\otimes w_j
\]
for some scalars $\lambda_{ij}\in\C$. It is important to remark that elements 
of $V\otimes W$ are not always of the form $v\otimes w$ for $v\in V$ and $w\in W$. 

The bilinearity of tensor products is crucial. 
For example,
\begin{align*}
    (v_1+v_3)\otimes (3w_1+w_2) 
    &=v_1\otimes (3w_1+w_2)+v_3\otimes (3w_1+w_2)\\
    &=v_1\otimes (3w_1)+v_1\otimes w_2+v_3\otimes (3w_1)+v_3\otimes w_2\\
    &=3v_1\otimes w_1+v_1\otimes w_2+3v_3\otimes w_1+v_3\otimes w_2.
\end{align*}

\begin{definition}
	Let $\rho\colon G\to\GL(V)$ and $\psi\colon G\to\GL(W)$ be representations, and 
    let $\{v_1,\dots,v_k\}$ be a basis of $V$ and $\{w_1,\dots,w_l\}$ a basis of $W$. 
    The \emph{tensor product} of $\rho$ and $\psi$ is the representation of $G$ given by 
	\begin{gather*}
	\rho\otimes\psi\colon G\to\GL(V\otimes W),
	\quad 
	g\mapsto (\rho\otimes\psi)_g,
	\shortintertext{where}
	(\rho\otimes\psi)_g\left (\sum_{i,j}\lambda_{ij} v_i\otimes w_j\right)=\sum_{i,j}\lambda_{ij} \rho_g(v_i)\otimes \psi_g(w_j).
	\end{gather*} 	
\end{definition}

\begin{exercise}
\label{xca:tensor}
Prove that the tensor product of representations 
is a representation.     
\end{exercise}



\section{Unitary representations}

\index{Inner product}
Since we are considering finite-dimensional vector spaces, our vector spaces admit 
an \emph{inner product}, that is a map $V\times V\to\C$, $(v,w)\mapsto\langle v,w\rangle$, 
that satisfies the following properties:
\begin{enumerate}
    \item $\langle y,x\rangle=\overline{\langle x,y\rangle}$ for all $x,y\in V$, where $\overline{a+bi}=a-bi$ denotes the complex conjugate.
    \item $\langle ax+by,z\rangle=a\langle x,z\rangle+b\langle y,z\rangle$ for all $x,y,z\in V$ and $a,b\in\C$.
    \item $\langle x,x\rangle\geq0$ for all $x\in V$, and 
      $\langle x,x\rangle=0$ if and only if $x=0$.
\end{enumerate}

\index{Inner product!usual}
For example, if $V=\C^n$, $v=\colvec{3}{v_1}{\vdots}{v_n}\in V$ and $w=\colvec{3}{w_1}{\vdots}{w_n}\in V$, 
then the \emph{usual inner product} between $v$ and $w$ 
is given by
\begin{equation}
    \label{eq:inner_product}
V\times V\to\C,\quad 
(v,w)\mapsto v_1\overline{w_1}+\cdots+v_n\overline{w_n}=v^T\overline{w},
\end{equation}
where 
\[
v^T=(v_1,\dots,v_n),\quad 
\overline{w}=\colvec{3}{\overline{w_1}}{\vdots}{\overline{w_n}}.
\]

\begin{definition}
    \index{Representation!unitary}
    Let $V$ be a finite-dimensional complex vector space with inner product $(v,w)\mapsto\langle v,w\rangle$. 
    A representation $\rho\colon G\to\GL(V)$ is \emph{unitary} if
    $\langle \rho_gv,\rho_gw\rangle=\langle v,w\rangle$ for all $g\in G$ and $v,w\in V$.
\end{definition}

\begin{exercise}
\index{Representation!left regular}
Let $G$ be a finite group and $V=\C[G]$ be the complex vector space with basis
$\{g:g\in G\}$. The \emph{left regular representation}
of $G$ is the representation
\[
L\colon G\to\GL(V),
\quad
g\mapsto L_g,
\]
where $L_g(h)=gh$. With respect to the inner product
\[
\left\langle\sum_{g\in G}\lambda_gg,\sum_{g\in G}\mu_gg\right\rangle=\sum_{g\in G}\lambda_g\overline{\mu_g},
\]
the representation $L$ is unitary. Prove it! 
\end{exercise}

\index{Orthogonal complement}
If $V$ is an inner-product vector space and $W$ is a subspace of $V$, 
\[
W^\perp = \{v\in V:\langle v,w\rangle=0\text{ for all $w\in W$}\}, 
\]
is called the 
\emph{orthogonal complement} of $W$. The following result is extremely important: 

\begin{proposition}
\label{pro:irr_or_dec}
Let $\rho\colon G\to\GL(V)$ be a unitary representation. Then $\rho$ is either
irreducible or decomposable.
\end{proposition}

\begin{proof}
	If $\rho$ is not irreducible, there exists an invariant subspace $W$ of $V$ such that
	$\{0\}\subsetneq W\subsetneq V$. Can we find an invariant complement? Yes, and this is because 
        our representation is unitary.  
	We know that $V=W\oplus W^{\perp}$ as complex vector spaces, where 
	$W^\perp\ne\{0\}$. So $\rho$ will be decomposable if we can prove that
	$W^\perp$ is an invariant subspace of $V$. 
	Let $v\in W^\perp$, $w\in W$ and $g\in G$. Then  
	\[
	\langle \rho_g(v),w\rangle=\langle v,\rho_g^{-1}(w)\rangle=0   
	\]
	since $\rho_g^{-1}(w)\in W$, as $W$ is an invariant subspace of $V$. This means that 
	$\rho_g(v)\in W^\perp$ and hence $\rho$ is decomposable. 
  %$V=W\oplus W^{\perp}$ is decomposable.  
\end{proof}

\section{Weyl's trick} 

\begin{theorem}[Weyl]
\index{Weyl's trick}
    Every representation $\rho\colon G\to\GL(V)$ 
    of a finite group is unitarizable, that is 
    $V$ admits an inner
    product such that
    \[
    \langle\rho_gv,\rho_gw\rangle=\langle v,w\rangle
    \]
    for all $g\in G$ and $v,w\in V$.
\end{theorem}

\begin{proof}
    Let $V\times V\to\C$, $(v,w)\mapsto v\cdot w$, be an\footnote{It does not
    matter which inner product we have chosen, we can safely use, for example,
  the usual inner product of \eqref{eq:inner_product}.} inner product on $V$. 
   
    A straightforward calculation shows that
    \[
    \langle v,w\rangle=\sum_{g\in G}(\rho_gv)\cdot (\rho_gw)
    \]
    is an inner product of $V$. For example, 
    \begin{align*}
    \langle v,w+w_1\rangle&=\sum_{g\in G}\rho_gv\cdot\rho_g(w+w_1)\\
    &=\sum_{g\in G}\rho_gv\cdot \left(\rho_gw+\rho_gw_1\right)\\
    &=\sum_{g\in G}\left(\rho_gv\cdot \rho_gw+\rho_gv\cdot\rho_gw_1\right)\\
    &=\sum_{g\in G}\rho_gv\cdot \rho_gw+\sum_{g\in G}\rho_gv\cdot \rho_gw_1\\
    &=\langle v,w\rangle+\langle v,w_1\rangle 
    \end{align*}
    holds for all $v,w,w_1\in V$. 
    
    Since
    \begin{align*}
    \langle\rho_gv,\rho_gw\rangle&=\sum_{h\in G}(\rho_h\rho_gv)\cdot(\rho_h\rho_gw)\\
    &=\sum_{h\in G}(\rho_{hg}v)\cdot(\rho_{hg}w)=\sum_{x\in G}(\rho_xv)\cdot(\rho_xw)=\langle v,w\rangle,
    \end{align*}
    the representation $\rho$ is unitary with respect to the newly 
    constructed inner product. 
\end{proof}

\begin{example}
    Let $A=\begin{pmatrix}
        -1 & -1\\
        1 & 0
    \end{pmatrix}$. Since $A^3=I$, we can use this matrix to 
    represent the (multiplicative) cyclic group 
    \[
    C_3=\langle g:g^3=1\rangle
    \]
    of order three. Note that the map $1\mapsto g$ yields an isomorphism 
    $\Z/3\simeq C_3$. 
        
    The map $C_3\to\GL_2(\C)$, $g\mapsto A$, is a group homomorphism. In particular, we have a representation $\rho\colon C_3\to\GL(V)$, where $V=\C^2$ with the usual inner product given by \eqref{eq:inner_product}, that is 
    \[
    v\cdot w=v^T\overline{w}
    \]
    for $v,w\in V$. 
    For example, the inner product between the vectors $\colvec{2}{1}{i}$ and $\colvec{2}{1}{2}$
    is 
    \[
    (1,i)\colvec{2}{1}{2}=1+2i.
    \]
    Since $A\in M_2(\R)$, we compute 
    \[
    \langle v,w\rangle=v^T\overline{w}+(Av)^T\overline{(Aw)}+(A^2v)^T\overline{(A^2w)}=v^T\left (I+A^TA+(A^2)^TA^2\right)\overline{w}
    \]
    and this is an invariant inner product on $V$. Note that 
    \[
        B=I+A^TA+(A^2)^TA^2=\begin{pmatrix}
            4&2\\
            2&4
        \end{pmatrix}.
    \]
\end{example}

%\label{rho_diagonalizable}
Weyl's trick has several interesting consequences:  

\begin{corollary}
\label{cor:consequences}
	Let $\rho\colon G\to\GL(V)$ be a representation of a finite group $G$. 
 The following properties
	hold:
	\begin{enumerate}
		\item $\rho$ is equivalent to a unitary representation.
		\item $\rho$ is either irreducible or decomposable.
		\item Each $\rho_g$ is diagonalizable. 
	\end{enumerate}
\end{corollary}

\begin{proof}
  The first assertion follows by choosing an orthonormal basis for the
  invariant inner product supplied by Weyl's unitary trick.  The second follows
  from 1), Proposition \ref{pro:irr_or_dec}, and Exercise
  \ref{xca:equivalence}.  Finally, with respect to an orthonormal basis for the
  invariant inner product, every $\rho_g$ is unitary and hence diagonalizable.
\end{proof}

\begin{exercise}
\label{xca:not_decomposable}
    If $G$ is an infinite group, it is no longer true that every non-zero representation
    is either irreducible or decomposable. Find an example.
\end{exercise}

For the previous exercise, consider the representation $\Z\to\GL_2(\C)$, 
$1\mapsto\begin{pmatrix}1&1\\0&1\end{pmatrix}$. 
\section{Maschke's theorem} 

Recall that, by convention, we only consider complex 
finite-dimensional representations of finite groups.

\begin{theorem}[Maschke]
\index{Maschke's theorem}
    Every representation of a finite group is completely reducible.
\end{theorem}

\begin{proof}
    Let $G$ be a finite group and $\rho\colon G\to\GL(V)$ be a representation of $G$. We proceed
    by induction on $\dim V$.
    If $\dim V=1$, the result is trivial, as degree-one representations are irreducible. Assume that
    the result holds for representations of degree $\leq n$. Suppose that $\rho$ has degree $n+1$. 
    If $\rho$ is irreducible, we are done. If not, use 
    Corollary~\ref{cor:consequences} to 
    write $V=S\oplus T$, where $S$ and $T$
    are non-zero invariant subspaces of $V$. Since $\dim S<\dim V$ and $\dim T<\dim V$, it follows from
    the inductive hypothesis that
    $\rho|_S$ and $\rho|_T$ are completely reducible. Since
    \[
    \rho\simeq\rho|_S\oplus\rho|_T,
    \]
    it follows that $\rho$ is completely reducible.
    %both $S$ and $T$ are spaces of completely reducible representations. 
    %Thus $\rho$ is completely reducible.
\end{proof}

\begin{example}
    Let $G=\Sym_3$ and $\rho\colon G\to\GL_3(\C)$ be the representation given by
    \[
    (12)\mapsto\begin{pmatrix}
    0&1&0\\
    1&0&0\\
    0&0&1
    \end{pmatrix},\quad
    (123)\mapsto\begin{pmatrix}
    0&0&1\\
    1&0&0\\
    0&1&0
    \end{pmatrix}.
    \]
    Then $\rho_g$ is unitary for all $g\in G$ (because $\rho_{(12)}$ and $\rho_{(123)}$ are both
    unitary). Moreover,
    \[
    S=\left\langle \begin{pmatrix}
    1\\1\\1
    \end{pmatrix}
    \right\rangle,
    \quad
    T=S^{\perp}=\left\langle
    \begin{pmatrix}
    -1\\1\\0
    \end{pmatrix},
    \begin{pmatrix}
    0\\-1\\1
    \end{pmatrix}
    \right\rangle,
    \]
    are invariant subspaces of $V=\C^3$, and the restricted representations on
    $S$ and $T$ are irreducible (see Exercise~\ref{xca:common_eigenvector}). A
    direct calculation shows that in the basis 
    \[
      \left\{\begin{pmatrix}
    1\\1\\1
    \end{pmatrix},
    \begin{pmatrix}
    -1\\1\\0
    \end{pmatrix},
    \begin{pmatrix}
    0\\-1\\1
    \end{pmatrix}
    \right\}
  \]
    the matrices $\rho_{(12)}$ and $\rho_{(123)}$ can be written as
    \[
    \rho_{(12)}=\begin{pmatrix}
        1&0&0\\
        0&-1&1\\
        0&0&1
    \end{pmatrix},
    \quad
    \rho_{(123)}=
    \begin{pmatrix}
        1&0&0\\
        0&0&-1\\
        0&1&-1
    \end{pmatrix}.
    \]
\end{example}

\section{Schur's lemma}

The following result is simple and crucial. 

\begin{lemma}[Schur]
\index{Schur's lemma}
    Let $\rho\colon G\to\GL(V)$ and $\psi\colon G\to\GL(W)$ be irreducible
    representations. If $T\colon V\to W$ is a non-zero invariant map, then $T$
    is bijective.  
\end{lemma}

\begin{proof}
    Since $T$ is non-zero and $\ker T$ is an invariant subspace of $V$, it
    follows that $\ker T=\{0\}$, as $\rho$ is irreducible. Thus $T$ is
    injective. Since $T(V)$ is a non-zero invariant subspace of $W$ and $\psi$
    is irreducible, it follows that $T$ is surjective. Therefore $T$ is
    bijective.  
\end{proof}

Two applications:

\begin{proposition}
\label{pro:Schur_consequence}
    If $\rho\colon G\to\GL(V)$ is an irreducible representation and $T\colon
    V\to V$ is invariant, then $T=\lambda\id$ for some $\lambda\in\C$. 
\end{proposition}

\begin{proof}
    Let $\lambda$ be an eigenvalue of $T$. Then $T-\lambda\id$ is invariant, as 
    \[
    (T-\lambda\id)\rho_g=T\rho_g-\lambda\rho_g=\rho_g(T-\lambda\id)
    \]
    for all $g\in G$ since $T$ is invariant. By definition, 
    $T-\lambda\id$ is not bijective. Thus $T-\lambda\id=0$ by Schur's lemma.
\end{proof}

\begin{proposition}
    Let $G$ be a finite abelian group. 
    If $\rho\colon G\to\GL(V)$ is an irreducible representation, then
    $\dim V=1$. 
\end{proposition}

\begin{proof}
    Let $h\in G$. Note that since $G$ is abelian, $T=\rho_h$ is invariant:
    \[
    T\rho_g=\rho_h\rho_g=\rho_{hg}=\rho_{gh}=\rho_g\rho_h=\rho_gT.
    \]
    By the previous proposition, there exists $\lambda_h\in\C$ such that
    $\rho_h=\lambda_h\id$. 
    %If $v\in V\setminus\{0\}$, then $V=\langle
    %v\rangle$. In fact, since every $\rho_h$ acts on $V$ by a scalar, $\langle
    %v\rangle$ is a non-zero invariant subspace of $V$ and $\rho$ is
    %irreducible, it follows that $V=\langle v\rangle$. Hence $\dim V=1$. 
    Let $v\in V\setminus\{0\}$. Since every $\rho_h$ acts on $V$ by a scalar,
    $\langle v\rangle$ is a non-zero invariant subspace of $V$. As $\rho$ is
    irreducible, $V=\langle v\rangle$ and hence $\dim V=1$.
\end{proof}

\begin{optional}
\section{Degree-one representations}

\index{Commutator subgroup}
The \emph{commutator subgroup} of a group $G$ is defined as the subgroup
$[G,G]$ generated by the commutators $[x,y]=xyx^{-1}y^{-1}$, that is  
 \[
        [G,G]=\langle[x,y]: x,y\in G\rangle.
\]

Routine calculations show that $[G,G]$ is always a normal subgroup of $G$. 
For example, the commutator subgroup of $\Z$ is the trivial subgroup $\{0\}$, 
as 
\[
[x,y]=x+y-x-y=0
\]
for all $x,y\in\Z$. As an exercise, one can also prove that 
$[\Sym_3,\Sym_3]=\{\id,(123),(132)\}$.

\begin{xca}
\label{xca:commutator}
Let $H$ be a normal subgroup of $G$. Prove that
$G/H$ is abelian if and only if $[G,G]\subseteq H$.
\end{xca}

\index{Abelianization}
The previous exercise demonstrates, in particular, that for any group 
$G$, the quotient $G/[G,G]$ is always an abelian group. 
This quotient is known as the \emph{abelianization} of 
$G$. 

\begin{xca}
Let $G$ be a finite group.
Prove that there is a bijection between degree-one representations of $G$ and
degree-one representations of $G/[G,G]$.
\end{xca}
\end{optional}


